\documentclass[12pt]{article}

% Pachete necesare pentru codificarea caracterelor și diacritice
\usepackage[utf8]{inputenc}
\usepackage[romanian]{babel}

% Pachete matematice avansate
\usepackage{amsmath}
\usepackage{amsfonts}
\usepackage{amssymb}
\usepackage{amsthm}

% Setarea marginilor paginii
\usepackage[margin=1in]{geometry}

\title{Seminar 2: Criptografie -- Rezolvare Teme}
\author{}
\date{} % Elimină complet data și spațiul ei de sub titlu

\begin{document}

\maketitle


\subsection*{Cerință}
\textit{Determinați inversul modular al lui $23$ modulo $25$.}


\subsection*{ Aplicarea Algoritmului lui Euclid Extins}

\subsubsection*{Pasul A: Algoritmul lui Euclid Standard (Împărțirile parțiale)}
Efectuăm împărțirile succesive cu rest pentru a confirma că numerele sunt prime între ele:
\begin{align*}
25 &= 23 \times 1 + 2 \\
23 &= 2 \times 11 + 1 \\
2 &= 1 \times 2 + 0
\end{align*}
Ultimul rest nenul este $1$, deci $\text{cmmdc}(25, 23) = 1$. Elementul este inversabil.

\subsubsection*{Pasul B: Substituția inversă (Determinarea coeficienților)}
Exprimăm resturile din relațiile de mai sus pentru a reconstitui identitatea lui Bézout:
\begin{align*}
1 &= 23 - 2 \times 11 \\
2 &= 25 - 23 \times 1
\end{align*}

Înlocuim expresia lui $2$ în prima ecuație:
\begin{align*}
1 &= 23 - (25 - 23 \times 1) \times 11 \\
1 &= 23 - 11 \times 25 + 11 \times 23 \\
1 &= 12 \times 23 + (-11) \times 25
\end{align*}

Am obținut coeficienții Bézout: $u = 12$ și $v = -11$.

\subsection*{Calculul Inversului Modular}
Trecând relația de mai sus în inelul claselor de resturi $\mathbb{Z}_{25}$, termenul alcătuit din multiplul lui $25$ devine congruent cu $0$:
\[
12 \times 23 + (-11) \times 25 = 1 \implies 12 \times 23 \equiv 1 \pmod{25}
\]

Prin urmare, inversul modular al lui $23$ modulo $25$ este:
\[
23^{-1} \equiv 12 \pmod{25}
\]

\subsection*{Verificare Aritmetică}
Pentru validarea rezultatului, înmulțim numărul inițial cu inversul determinat:
\[
23 \times 12 = 276
\]
Efectuăm împărțirea cu rest la modulul $25$:
\[
276 = 25 \times 11 + 1 \implies 276 \equiv 1 \pmod{25}
\]
Relația este verificată și complet demonstrată.

\end{document}